\documentclass[12pt,a4paper]{article}
\usepackage[utf8]{inputenc}
\usepackage{amsmath,amssymb,amsfonts,amsthm}
\usepackage{geometry}
\usepackage{hyperref}
\usepackage{xcolor}

\geometry{a4paper, margin=1in}

\title{\textbf{\Huge Thử nghiệm Live Preview \LaTeX\ trên Neovim}}
\author{\textbf{Waanne} -- NixOS Environment}
\date{\today}

\newtheorem{theorem}{Định lý}

\begin{document}

\maketitle

\section{Giới thiệu tính năng Live Preview}
\subsection{Các phím tắt quan trọng trong Neovim}
\begin{itemize}
	\item \texttt{<leader>lk} (Space + l + k): Dừng tiến trình biên dịch nền (\textit{Stop compiler}).
	\item \texttt{<leader>lc} (Space + l + c): Dọn dẹp các file rác sinh ra trong quá trình build (\textit{Clean auxiliary files}).
\end{itemize}

\section{Công thức Toán học \& Định lý mẫu}

Dưới đây là một số công thức toán học đẹp mắt để kiểm tra khả năng hiển thị của \LaTeX:

\begin{theorem}[Đẳng thức Euler kỳ diệu]
	Đẳng thức liên kết 5 hằng số cơ bản nhất của toán học:
	\begin{equation}
		e^{i\pi} + 1 = 0
	\end{equation}
\end{theorem}

\subsection{Tích phân Gaussian và Ma trận}
Tích phân suy rộng nổi tiếng của hàm Gaussian:
\begin{equation}
	\int_{-\infty}^{+\infty} e^{-x^2} \, dx = \sqrt{\pi}
\end{equation}

Ma trận xoay trong không gian 2D với góc quay $\theta$:
\begin{equation}
	R(\theta) = \begin{pmatrix}
		\cos\theta & -\sin\theta \\
		\sin\theta & \cos\theta
	\end{pmatrix}
\end{equation}

\subsection{Thử nghiệm chỉnh sửa trực tiếp}
Hãy thử sửa dòng chữ này hoặc thêm công thức mới:
\begin{equation}
	f(x) = \sum_{n=0}^{\infty} \frac{f^{(n)}(a)}{n!} (x - a)^n
\end{equation}
Sau khi sửa, bạn chỉ cần bấm \textbf{\texttt{:w}} và quan sát cửa sổ Zathura tự động làm mới!

\end{document}
